% =============================================================================
% preambulo.tex · Paquetes, formato y entornos de la tesis maestra
% Formato provisional (decisión B, 20-10-2026): clase book con los márgenes y el
% espaciado de tesis/plantilla_uni.docx. Debe revisarse cuando se disponga de las
% «Pautas para redactar el Trabajo de Tesis de Maestría» de la UNI.
% =============================================================================

% ---- Página: A4, 3 cm a la izquierda y 2,5 cm en los demás lados ----------------
% La plantilla Word no fija el tamaño de papel; se adopta A4, el estándar en el Perú.
\usepackage[a4paper, left=3cm, right=2.5cm, top=2.5cm, bottom=2.5cm,
            footskip=1.25cm, headheight=15pt]{geometry}

% ---- Codificación y tipografía ----------------------------------------------
% El preámbulo no usa fontspec: en este entorno TeX Live no incluye luaotfload (paquete
% texlive-luatex), sin el cual fontspec no funciona. Se usan fuentes Type 1 con
% codificación T1, que compilan igual en cualquier instalación de TeX Live.
% Texto: Nimbus Sans (paquete helvet), clon de Helvetica; Arial reproduce los anchos de
% Helvetica, de modo que es el equivalente métrico de Arial 12. Matemáticas: Computer
% Modern (amsfonts), porque no hay una fuente matemática sin serifas instalada.
\usepackage[T1]{fontenc}
\usepackage{textcomp}
\usepackage{helvet}
\renewcommand{\familydefault}{\sfdefault}
\renewcommand{\rmdefault}{phv}       % evita sustituciones: no hay Latin Modern T1 instalada
\renewcommand{\ttdefault}{pcr}       % monoespaciada: Courier (Type 1; las EC solo existen en mapa de bits)
% Sin luaotfload, LuaLaTeX pasa el código Unicode a la fuente T1, que coincide con Latin-1
% en las letras acentuadas pero no en los signos siguientes; se definen explícitamente.
\usepackage{newunicodechar}
\newunicodechar{«}{\guillemotleft}
\newunicodechar{»}{\guillemotright}
\newunicodechar{¿}{\textquestiondown}
\newunicodechar{¡}{\textexclamdown}
\newunicodechar{—}{\textemdash}
\newunicodechar{–}{\textendash}
\newunicodechar{“}{\textquotedblleft}
\newunicodechar{”}{\textquotedblright}
\newunicodechar{‘}{\textquoteleft}
\newunicodechar{’}{\textquoteright}
\newunicodechar{·}{\textperiodcentered}
\newunicodechar{ń}{\'n}                % Caliński, en el .bib
\usepackage{amsmath, amssymb, amsthm, mathtools, bm}
\usepackage{icomma}                  % 0,7 en modo matemático sin espacio tras la coma

% ---- Idioma -----------------------------------------------------------------
% es-tabla: «Tabla» en lugar de «Cuadro»; es-noshorthands evita conflictos de «"» y «~»;
% es-nodecimaldot conserva el punto en modo matemático (la coma decimal la da icomma).
\usepackage[spanish, es-tabla, es-noshorthands, es-nodecimaldot]{babel}
\usepackage{csquotes}
\DeclareQuoteStyle{spanish}{\guillemotleft}{\guillemotright}{\textquotedblleft}{\textquotedblright}

% ---- Interlineado 1,5 (equivalente al w:line="360" de la plantilla Word) -------
\usepackage{setspace}
\onehalfspacing
\setlength{\parindent}{0pt}
\setlength{\parskip}{6pt plus 2pt minus 1pt}
\setlength{\emergencystretch}{3em}   % evita líneas desbordadas en párrafos con términos largos

% ---- Encabezados y pie: número de página centrado, como la plantilla ----------
\usepackage{fancyhdr}
\pagestyle{fancy}
\fancyhf{}
\fancyhead[R]{\small\nouppercase{\leftmark}}
\fancyfoot[C]{\thepage}
\renewcommand{\headrulewidth}{0pt}
\fancypagestyle{plain}{\fancyhf{}\fancyfoot[C]{\thepage}}

% ---- Títulos ----------------------------------------------------------------
\usepackage{titlesec}
\titleformat{\chapter}[display]{\normalfont\bfseries\Large}
  {\MakeUppercase{\chaptertitlename}\ \thechapter}{0.6em}{\LARGE}
\titlespacing*{\chapter}{0pt}{0pt}{1.5em}
\titleformat{\section}{\normalfont\bfseries\large}{\thesection}{0.8em}{}
\titleformat{\subsection}{\normalfont\bfseries\normalsize}{\thesubsection}{0.8em}{}
\titleformat{\subsubsection}{\normalfont\bfseries\itshape\normalsize}{\thesubsubsection}{0.8em}{}
\setcounter{secnumdepth}{3}
\setcounter{tocdepth}{2}

% ---- Tablas, figuras y listas ------------------------------------------------
\usepackage{booktabs, longtable, array}
\usepackage{graphicx}
\graphicspath{{../reports/figures/}{../reports/figures/niveles/}}
\usepackage[font=small, labelfont=bf]{caption}
\usepackage{enumitem}
\setlist{itemsep=2pt, topsep=4pt}
\usepackage{xcolor}

% ---- Bibliografía: APA 7 con biblatex-apa y el .bib compartido con el informe --
\usepackage[style=apa, backend=biber]{biblatex}
\addbibresource{../tesis/referencias.bib}

% ---- Enlaces y referencias cruzadas (hyperref antes de cleveref) --------------
\usepackage[hidelinks, unicode, pdfusetitle]{hyperref}
\usepackage[spanish, noabbrev]{cleveref}

% ---- Entornos de amsthm, numerados por capítulo -----------------------------
\theoremstyle{plain}
\newtheorem{teorema}{Teorema}[chapter]
\newtheorem{proposicion}[teorema]{Proposición}
\newtheorem{lema}[teorema]{Lema}
\newtheorem{corolario}[teorema]{Corolario}
\theoremstyle{definition}
\newtheorem{definicion}[teorema]{Definición}
\newtheorem{supuesto}[teorema]{Supuesto}
\newtheorem{algoritmo}[teorema]{Algoritmo}
\theoremstyle{remark}
\newtheorem{observacion}[teorema]{Observación}
\renewcommand{\qedsymbol}{$\blacksquare$}
\addto\captionsspanish{\renewcommand{\proofname}{Demostración}}

\crefname{teorema}{teorema}{teoremas}
\Crefname{teorema}{Teorema}{Teoremas}
\crefname{proposicion}{proposición}{proposiciones}
\Crefname{proposicion}{Proposición}{Proposiciones}
\crefname{lema}{lema}{lemas}
\Crefname{lema}{Lema}{Lemas}
\crefname{corolario}{corolario}{corolarios}
\Crefname{corolario}{Corolario}{Corolarios}
\crefname{definicion}{definición}{definiciones}
\Crefname{definicion}{Definición}{Definiciones}
\crefname{supuesto}{supuesto}{supuestos}
\Crefname{supuesto}{Supuesto}{Supuestos}
\crefname{algoritmo}{algoritmo}{algoritmos}
\Crefname{algoritmo}{Algoritmo}{Algoritmos}
\crefname{observacion}{observación}{observaciones}
\Crefname{observacion}{Observación}{Observaciones}
\crefname{equation}{ecuación}{ecuaciones}
\Crefname{equation}{Ecuación}{Ecuaciones}
\creflabelformat{equation}{#2(#1)#3}
\crefname{chapter}{capítulo}{capítulos}
\Crefname{chapter}{Capítulo}{Capítulos}
\crefname{section}{sección}{secciones}
\Crefname{section}{Sección}{Secciones}
\crefname{subsection}{subsección}{subsecciones}
\Crefname{subsection}{Subsección}{Subsecciones}
\crefname{enumi}{objetivo}{objetivos}
\Crefname{enumi}{Objetivo}{Objetivos}
\numberwithin{equation}{chapter}

% ---- Marcas editoriales -----------------------------------------------------
% \pendiente{qué falta}: obra sin ficha en papers/ o cifra aún no generada. Se resalta en
%   el PDF y cada uso se lista en la nota del redactor y en tesis/registro_verificacion.md.
% \verificar{motivo}: obra con ficha y entrada en el .bib, pero con un dato por verificar
%   (equivale a la marca [VERIFICAR: motivo] del informe). Va como nota al pie.
\newcommand{\pendiente}[1]{%
  \textcolor{red!75!black}{\textbf{[PENDIENTE:}~\textit{#1}\textbf{]}}}
\newcommand{\verificar}[1]{%
  \footnote{\textcolor{orange!70!black}{\textbf{VERIFICAR:}}~#1}}
